\documentclass[../main.tex]{subfiles}
\begin{document}

\chapter{Storage}
\lessoninfo{
  video={Lesson 16, videos 1--14},
  textbook={H\&P \cite{hennessy2017} §2.2, §D.1--D.2},
  keywords={magnetic disk, platter, cylinder, sector, seek time, rotational latency, optical disk, magnetic tape, sequential access, flash memory, solid-state disk, hybrid disk, I/O bus, mezzanine bus},
}

Lesson 15 showed that main memory is built from DRAM, which loses its
contents as soon as power is removed.  Programs and data that must survive
have to be kept on storage devices, which also hold the pages that virtual
memory (Lesson 13) cannot fit into physical memory.  This lesson describes
the main storage technologies---magnetic disks, optical disks, magnetic tape
and semiconductor storage based on DRAM or flash memory---with emphasis on
what determines their access time and how it evolves, and it shows how
storage and other I/O devices are connected to the processor.

\section{Storage}
\label{sec:l16-storage}

Storage keeps data \emph{persistently}: unlike main memory, a storage device
retains its contents when the computer is switched off.\source{Lesson 16,
videos 1--2; H\&P §D.1.}  It serves two purposes.
\begin{description}
  \item[Files] Programs, the data they work on and the settings of the
    system and its applications are kept in files, which must survive from
    one session to the next.
  \item[Virtual memory] The pages that do not fit into physical memory are
    kept on storage, and a page fault reads such a page back into memory
    (Lesson 13).
\end{description}

Two properties of a storage device matter to the architect.  Its
\emph{performance} is described by the two metrics of Lesson 2: throughput,
the amount of data or the number of requests handled per second, and
latency, the time a single access takes.  Storage throughput improves over
time, but more slowly than processor speed; storage latency improves even
more slowly, more slowly than the latency of DRAM, so the gap between the
processor and storage is wider still than the memory wall of Lesson 1.\margin{Read
latency: DRAM 60--\SI{90}{\nano\second}, a flash SSD about
\SI{100}{\micro\second}, a disk about \SI{10}{\milli\second} (7-cpu.com;
vendor data sheets).}  Its
\emph{reliability} matters at least as much: storage is where data is
expected to be safe, and a lost file is often worse than a crashed program.
Reliability is the subject of Lesson 17.

\section{Magnetic disks}
\label{sec:l16-disk}

A \term{magnetic disk}[磁気ディスク]\index{hard disk|see{magnetic disk}}
(hard disk) records bits as the direction of magnetization of tiny regions of
a magnetic coating.\source{Lesson 16, video 3; H\&P §D.1.}  It is built
from rotating disks covered with this coating and from heads that move over
them (\cref{fig:l16-disk}).

\begin{figure}[htbp]
  \centering
  % Structure of a magnetic disk.  (a) Edge-on view: platters on a spindle,
% one head per surface on arms that form the head assembly; the dashed
% lines mark one cylinder.  (b) Top view of one surface: concentric tracks
% divided into sectors; one track and one sector are highlighted.
\begin{tikzpicture}[x=1cm,y=1cm, font=\footnotesize\sffamily, >={Stealth[length=1.6mm]},
  lab/.style={font=\scriptsize\sffamily, text=ln-muted},
  pin/.style={draw=ln-muted, thin}]
  % ================================================= (a) edge-on view
  \node[anchor=west] at (-0.2,3.35) {\iflangja{(a) 側面}{(a) side view}};
  % spindle
  \draw[fill=ln-muted!35] (1.45,-0.15) rectangle (1.75,2.45);
  % platters
  \foreach \y in {0.4,1.2,2.0} {
    \draw[fill=ln-accent!30] (0.1,\y-0.05) rectangle (3.1,\y+0.05);
  }
  % one cylinder: all tracks at the same radius
  \draw[ln-accent, dashed, thick] (2.35,0.15) -- (2.35,2.25);
  \draw[ln-accent, dashed, thick] (0.85,0.15) -- (0.85,2.25);
  % head assembly: actuator and arms, one head per surface
  \draw[fill=ln-box] (4.3,0.12) rectangle (4.55,2.28);
  \foreach \y in {0.4,1.2,2.0} {
    \foreach \s in {1,-1} {
      \draw[very thick, ln-muted] (4.3,\y+\s*0.15) -- (2.5,\y+\s*0.15);
      \fill[ln-accent] (2.22,\y+\s*0.07) rectangle (2.5,\y+\s*0.19);
    }
  }
  % rotation
  \draw[->, thick] (1.05,2.62) arc[start angle=190, end angle=350, x radius=0.55, y radius=0.14];
  % labels
  \node[lab, anchor=south] at (1.6,2.75) {\iflangja{スピンドル}{spindle}};
  \draw[pin] (0.35,0.35) -- (0.35,-0.25);
  \node[lab, anchor=north] at (0.35,-0.22) {\iflangja{プラッタ}{platter}};
  \draw[pin] (2.35,0.15) -- (2.35,-0.25);
  \node[lab, anchor=north] at (2.35,-0.22) {\iflangja{シリンダ}{cylinder}};
  \draw[pin] (2.42,2.2) -- (3.2,2.74);
  \node[lab, anchor=south] at (3.3,2.72) {\iflangja{ヘッド}{head}};
  \node[lab, anchor=north] at (4.43,0.05) {\iflangja{ヘッドアセンブリ}{head assembly}};
  % ================================================= (b) top view
  \begin{scope}[xshift=6.3cm]
    \node[anchor=west] at (-0.2,3.35) {\iflangja{(b) 1 面の上面}{(b) top view of one surface}};
    \coordinate (c) at (2.0,1.2);
    \draw[fill=ln-box] (c) circle[radius=1.5];
    \foreach \r in {0.5,0.65,...,1.4} {
      \draw[ln-rule] (c) circle[radius=\r];
    }
    \foreach \a in {0,30,...,330} {
      \draw[ln-rule] ($(c)+(\a:0.5)$) -- ($(c)+(\a:1.4)$);
    }
    % one track and one sector
    \draw[ln-accent, thick] (c) circle[radius=0.95];
    \fill[ln-accent] ($(c)+(30:1.02)$) arc[start angle=30, end angle=60, radius=1.02]
      -- ($(c)+(60:0.88)$) arc[start angle=60, end angle=30, radius=0.88] -- cycle;
    % spindle
    \draw[fill=ln-muted!35] (c) circle[radius=0.3];
    % head on its arm
    \draw[very thick, ln-muted] (3.75,-0.05) -- ($(c)+(-12:0.95)$);
    \fill[ln-muted] (3.75,-0.05) circle[radius=0.08];
    \fill[ln-accent] ($(c)+(-12:0.95)$) circle[radius=0.07];
    % rotation
    \draw[->, thick] ($(c)+(105:1.65)$) arc[start angle=105, end angle=165, radius=1.65];
    % labels
    \draw[pin] ($(c)+(45:0.95)$) -- (3.45,2.7);
    \node[lab, anchor=south] at (3.45,2.67) {\iflangja{セクタ}{sector}};
    \draw[pin] ($(c)+(215:0.95)$) -- (0.1,-0.2);
    \node[lab, anchor=north] at (0.1,-0.17) {\iflangja{トラック}{track}};
    \draw[pin] ($(c)+(-12:0.95)$) -- (3.55,1.45);
    \node[lab, anchor=west] at (3.52,1.5) {\iflangja{ヘッド}{head}};
  \end{scope}
\end{tikzpicture}

  \caption{Structure of a magnetic disk.  (a)~Platters on a spindle, with one
    head per surface; the arms of all heads form the head assembly, and the
    dashed lines mark one cylinder.  (b)~One surface, divided into
    concentric tracks and each track into sectors; one track and one of its
    sectors are highlighted.}
  \label{fig:l16-disk}
\end{figure}

The coating covers both surfaces of each \term{platter}[プラッタ], a rigid
circular disk.  Several platters are stacked on a common
\term{spindle}[スピンドル], which a motor turns at a constant speed, so that
all platters rotate together (\cref{fig:l16-disk}a).

Each surface is read and written by its own
\term{disk head}[磁気ヘッド], which is mounted at the end of an arm and
floats a tiny distance above the surface without touching it; a head that
touches the rotating surface damages it.  All arms are attached to one
actuator and move together as the \emph{head assembly}, which can place the
heads anywhere between the outer edge of the platters and the spindle.

With the head assembly held in one position, each head passes over a circle
of its surface as the platter rotates.  This circle is a
\term{track}[トラック], and a surface holds many concentric tracks
(\cref{fig:l16-disk}b).  Because the heads move together, they are always
positioned over tracks at the same distance from the spindle, one on each
surface.  These tracks together form a \term{cylinder}[シリンダ]: moving the
head assembly selects a cylinder, and selecting one of the heads selects a
surface, and therefore a single track, within it.

Along the track the data is divided into blocks of fixed size.  Such a block,
a \term{sector}[セクタ], is the smallest unit that the disk reads or writes.
A sector begins with a short preamble, which lets the disk recognize the
start of the sector and synchronize with its bits; the data follows, and the
sector ends with an error-correcting code with which the disk checks, and
where possible corrects, the data it reads (Lesson 17).

The capacity of a disk follows directly from this organization:
\begin{equation}
  \text{capacity} \;=\; \#\text{platters} \times 2
    \times \frac{\text{tracks}}{\text{surface}}
    \times \frac{\text{sectors}}{\text{track}}
    \times \frac{\text{bytes}}{\text{sector}} ,
  \label{eq:l16-capacity}
\end{equation}
where the factor 2 counts the two surfaces of a platter and the number of
tracks per surface equals the number of cylinders.\margin{Real disks place
more sectors on the longer outer tracks; sectors per track is then an
average.}

\begin{example}
  \label{ex:l16-capacity}
  A disk has two platters with \num{100000} tracks per surface, on average
  \num{2000} sectors per track, and \num{4096} bytes per sector.  By
  \cref{eq:l16-capacity} it holds $2\times 2\times \num{100000}\times
  \num{2000} = \num{8e8}$ sectors, or $\num{8e8}\times \num{4096} \approx
  \num{3.3e12}$ bytes, i.e.\ about 3.3~TB.\caution{Disk capacities are
  decimal: 1~TB is $10^{12}$ bytes, not $2^{40}$, so a 4~TB drive holds only
  3.64~TiB.}
\end{example}

\section{Disk access time}
\label{sec:l16-access}

To read or write a sector, the disk has to bring the right head over the
right place of the rotating surface, and then the data has to travel to
main memory.\source{Lesson 16, videos 4--5; H\&P §D.1.}  For a disk whose
platters are already spinning, the time of one access consists of five
parts.
\begin{description}
  \item[Seek time] The \term{seek time}[シーク時間] is the time the head
    assembly needs to move the heads to the cylinder that contains the
    sector.  It grows with the distance the heads have to travel.
  \item[Rotational latency] The \term{rotational latency}[回転待ち時間] is
    the time from the arrival of the head at the track until the start of
    the sector rotates under the head.
  \item[Data read time] The sector passes under the head at the speed of
    rotation, and it is read (or written) during this time; an access to
    several consecutive sectors takes correspondingly longer.
  \item[Controller time] The disk controller processes the request, for
    example checking the error-correcting code of each sector read.
  \item[I/O bus time] The data is transferred between the disk and main
    memory over the I/O bus (\cref{sec:l16-io}).
\end{description}
\begin{equation}
  t_{\text{access}} \;=\; t_{\text{seek}} + t_{\text{rot}} + t_{\text{read}}
    + t_{\text{controller}} + t_{\text{bus}} .
  \label{eq:l16-access}
\end{equation}
If the disk is still busy with earlier requests, a new request also waits in
a queue before its own access begins.  This queueing delay, the time until
the earlier accesses have finished and the heads can move again, is often
significant.

The two rotation-dependent terms follow from the rotation speed, usually
given in revolutions per minute (RPM).  One revolution takes $60/\text{RPM}$
seconds.\tip{In milliseconds a revolution takes $\num{60000}/\text{RPM}$: at
\num{7200}~RPM it is \SI{8.33}{\milli\second}, so the average rotational
latency is about \SI{4.2}{\milli\second}.}  The start of the sector may be
anywhere on the track when the head arrives, so the rotational latency lies
between zero and one revolution and is half a revolution on average.  Reading
$n$ consecutive sectors takes the fraction $n/(\text{sectors per track})$ of a
revolution:
\begin{equation}
  \overline{t_{\text{rot}}} \;=\; \frac{1}{2}\cdot\frac{60}{\text{RPM}} ,
  \qquad
  t_{\text{read}} \;=\; \frac{n}{\text{sectors per track}}\cdot
    \frac{60}{\text{RPM}} .
  \label{eq:l16-rot}
\end{equation}

\begin{example}
  \label{ex:l16-access}
  A disk rotates at \num{10000}~RPM and has an average seek time of
  \SI{4}{\milli\second}.  A request reads 100 consecutive sectors of a track
  that holds 500 sectors; the controller needs \SI{0.3}{\milli\second} and
  the transfer over the I/O bus \SI{0.5}{\milli\second}.  One revolution
  takes $\num{60000}/\num{10000} = \SI{6}{\milli\second}$, so by
  \cref{eq:l16-rot} the average rotational latency is
  \SI{3}{\milli\second} and the data read time is $100/500 \times
  \SI{6}{\milli\second} = \SI{1.2}{\milli\second}$.  By
  \cref{eq:l16-access} the access takes on average $4 + 3 + 1.2 + 0.3 + 0.5
  = \SI{9}{\milli\second}$ (\cref{fig:l16-access-time}).  Had the 100
  sectors been scattered over the disk, each would have needed its own seek
  and rotational latency, and the request would have taken more than
  \SI{700}{\milli\second}.\margin{A drive serves one access at a time:
  \SI{9}{\milli\second} is about 110 requests per second.  Vendors quote
  this as IOPS; a flash SSD reaches \num{100000} and more.}
\end{example}

\begin{figure}[htbp]
  \centering
  % Components of one disk access, drawn to scale for the numbers of
% Example ex:l16-access (seek 4 ms, rotational latency 3 ms, data read
% 1.2 ms, controller 0.3 ms, I/O bus 0.5 ms).  1 ms = 1.2 cm.
\begin{tikzpicture}[x=1.2cm,y=1cm, font=\footnotesize\sffamily,
  lab/.style={font=\scriptsize\sffamily, text=ln-muted},
  pin/.style={draw=ln-muted, thin},
  seg/.style={draw, minimum height=0.55cm}]
  % segments (x in ms)
  \draw[fill=ln-accent!40] (0,0) rectangle (4,0.55);
  \draw[fill=ln-accent!20] (4,0) rectangle (7,0.55);
  \draw[fill=ln-accent!10] (7,0) rectangle (8.2,0.55);
  \draw[fill=ln-box] (8.2,0) rectangle (8.5,0.55);
  \draw[fill=white] (8.5,0) rectangle (9,0.55);
  \node at (2,0.275) {\iflangja{シーク時間}{seek time}};
  \node at (5.5,0.275) {\iflangja{回転待ち時間}{rotational latency}};
  \node[lab, anchor=south] at (7.6,0.6) {\iflangja{データ読み出し}{data read}};
  \draw[pin] (8.35,0.55) -- (8.35,1.0);
  \node[lab, anchor=south] at (8.35,0.97) {\iflangja{コントローラ}{controller}};
  \draw[pin] (8.75,0.55) -- (8.95,0.72);
  \node[lab, anchor=south west] at (8.85,0.62) {\iflangja{I/O バス}{I/O bus}};
  % time axis
  \draw (0,-0.1) -- (9,-0.1);
  \foreach \t in {0,...,9} {
    \draw (\t,-0.1) -- (\t,-0.18);
    \node[lab, anchor=north] at (\t,-0.16) {\t};
  }
  \node[lab, anchor=north west] at (9.1,-0.16) {ms};
  % mechanical vs electronic parts
  \draw[ln-muted] (0,-0.6) -- (0,-0.68) -- (8.15,-0.68) -- (8.15,-0.6);
  \node[lab, anchor=north] at (4.1,-0.7) {\iflangja{機械的な動作（ヘッドとプラッタ）}{mechanical (head and platter movement)}};
  \draw[ln-muted] (8.25,-0.6) -- (8.25,-0.68) -- (9,-0.68) -- (9,-0.6);
  \node[lab, anchor=north] at (8.61,-0.7) {\iflangja{電子的}{electronic}};
\end{tikzpicture}

  \caption{The access of \cref{ex:l16-access}, drawn to scale.  Mechanical
    movement accounts for \SI{8.2}{\milli\second} of the
    \SI{9}{\milli\second}.}
  \label{fig:l16-access-time}
\end{figure}

\tip{Read \SI{1}{\milli\second} as one million cycles per gigahertz: the
\SI{9}{\milli\second} above is 27 million cycles on a \SI{3}{\giga\hertz}
processor.}

\begin{keypoint}
  A disk access is dominated by mechanical movement: moving the head
  assembly and waiting for the platter to turn.  It takes milliseconds,
  i.e.\ millions of clock cycles of a processor, which is why a page fault
  (Lesson 13) is so expensive.
\end{keypoint}

\section{Trends for magnetic disks}
\label{sec:l16-trends}

The components of a magnetic disk improve at very different
rates.\source{Lesson 16, video 6; H\&P §D.2.}
\begin{description}
  \item[Capacity] depends on how densely bits are recorded on the surface.
    For many years it doubled roughly every one to two years; this growth
    has slowed considerably since about 2010.
  \item[Seek time] is typically \SIrange{5}{10}{\milli\second} and improves
    only very slowly.\caution{A quoted seek time is an average over random
    seeks; a track-to-track seek takes under \SI{1}{\milli\second}, a
    full-stroke seek about twice the average.}  What improvement there is
    comes mainly from smaller
    platter diameters, which shorten the distance the heads travel.
  \item[Rotation speed] rose slowly, from about \num{5000}~RPM to
    \num{10000}--\num{15000}~RPM for the fastest drives; most drives spin at
    \num{5400}--\num{7200}~RPM.  A faster rotation requires stronger
    materials and produces more noise, heat and power consumption.
  \item[Controller and bus time] are electronic and improve at a good rate,
    but they are a small part of the access time (\cref{fig:l16-access-time}).
\end{description}

The latency of a disk access is therefore set by mechanics and materials
rather than by the progress of transistor technology that Moore's law
describes (Lesson 1).  It improves far more slowly than processor speed, even
though capacity grew quickly for many years.\sidenote{The magnetic disk is
the oldest
component of a computer still in use.  IBM's 350 disk storage unit (1956)
held about 5~MB on fifty 24-inch platters turning at \num{1200}~RPM, and its
single head needed some \SI{0.6}{\second} to reach a track.  Capacity has
since grown by six orders of magnitude, the access time by less than two.}
Higher recording density does
shorten the data read time of a given amount of data, since more bytes pass
under the head in each revolution, but it does not affect the seek time or
the rotational latency.

\section{Optical disks}
\label{sec:l16-optical}

An \term{optical disk}[光ディスク] is a single rotating disk like one platter
of a magnetic disk, with data along a track, but it records bits as marks that
change how the surface reflects light, and it reads them with a laser
instead of a magnetic head.\source{Lesson 16, video 7.}  The laser works at a
distance from the disk and focuses through a transparent protective layer,
so dust and fingerprints on the surface are far less harmful than on a
magnetic platter, which must be sealed inside its drive.  An optical disk can
therefore be removed from its drive and carried to another one: it is
\emph{portable}.\margin{On CDs, DVDs and Blu-ray discs the data lies along a
single spiral track.}

Portability has a price.  A disk written in one drive must be readable in
every other drive, so the recording format has to be fixed by an industry
standard: CD, then DVD, then Blu-ray.\margin{CD (1982) holds 700~MB, DVD
(1996) 4.7~GB and Blu-ray (2006) 25~GB per layer: roughly a decade and a
factor of six apart.}  A denser recording technology becomes
useful only after a new standard has been agreed on and drives, media and
content for it have become widespread, which takes years.  Optical disks
therefore improve in large, infrequent steps.  The platters of a magnetic
disk never leave their drive; a manufacturer can adopt a better recording
technology in its next product as long as the drive still connects through
the same standard interface.

\section{Magnetic tape}
\label{sec:l16-tape}

A \term{magnetic tape}[磁気テープ] records data on a long, thin strip
coated with magnetic material, which is kept wound up in a cartridge and
pulled past the read/write head of the drive.\source{Lesson 16, videos
8--9.}  A tape has a large capacity, and its cartridges are replaceable: one
drive can read and write any number of cheap cartridges.  Its drawback is
the way data is reached.

\begin{definition}[Sequential access]
  A medium provides \term{sequential access}[順次アクセス] if data can be
  reached only in order: to access a location, the medium has to be moved
  past all the data between the current position and the requested one, so
  the access time grows with that distance.
\end{definition}

A disk, by contrast, reaches any sector directly, without passing over the
data in between: one seek and at most one revolution take at most some tens
of milliseconds, wherever the sector lies (\cref{ex:l16-tape}).  Its access
time still depends on the seek distance, so a disk provides random access
(Lesson 15) only approximately, but far more closely than a tape.  Sequential
access suits one use particularly well: \emph{backup}, where large amounts of
data are written in one pass and read back rarely, and then usually in one
pass as well.\margin{An LTO-9 cartridge (2021) holds 18~TB and streams at
400~MB/s, so reading it end to end takes about 12 hours (LTO consortium).}
It is unsuitable for virtual memory, whose page faults read
pages from arbitrary locations.  Tape is therefore used as backup storage
rather than as the storage from which programs and files are used.

\begin{example}
  \label{ex:l16-tape}
  A tape of \SI{1000}{\metre} is wound at
  \SI[per-mode=symbol]{5}{\metre\per\second}.  With the tape at its
  beginning, data in the middle is reached after
  \SI{100}{\second} and data at the end after \SI{200}{\second}.  A disk
  whose longest seek takes \SI{15}{\milli\second} and whose revolution takes
  \SI{6}{\milli\second} reaches any sector within about
  \SI{21}{\milli\second}, about \num{9500} times faster than the tape
  reaches its end.  For reading a whole backup from start to finish, the
  winding time does not matter.
\end{example}

Tape is gradually being displaced.  Its production volume is small, so its
cost per gigabyte falls more slowly than that of magnetic disks, which are
made in huge numbers; disks, for example external disks attached over USB,
have become cheap enough to take over much of the backup role.

\section{Solid-state disks}
\label{sec:l16-ssd}

Semiconductor memory has no moving parts, and an access to it takes
nanoseconds to a fraction of a millisecond instead of the milliseconds of a
disk.\source{Lesson 16, video 10; H\&P §2.2, §D.1.}  Its cost per gigabyte,
however, is much higher.  DRAM costs on the order of a hundred times more per
gigabyte than a magnetic disk, while its latency is about \num{100000} times
lower (tens of nanoseconds against milliseconds).  Storage built from
semiconductor memory is used where speed justifies this cost.

\begin{definition}[Solid-state disk]
  A \term{solid-state disk}[SSD]\index{SSD|see{solid-state disk}} (SSD) is a
  storage device built from semiconductor memory, without moving parts, that
  is used like a disk: it stores blocks of data and connects through a disk
  interface such as SATA or directly over PCI Express.
\end{definition}

An SSD can be built from either of two kinds of memory.
\begin{description}
  \item[DRAM with a battery] DRAM is volatile, so a battery keeps it powered
    when the computer is off.  Such an SSD is extremely fast and extremely
    expensive per gigabyte, and its data is safe only as long as the battery
    lasts, which makes it unsuitable for keeping data over long periods.
  \item[Flash memory] Most SSDs use \term{flash memory}[フラッシュメモリ],
    a semiconductor memory that keeps its contents without power: its bits
    are stored as charge trapped in transistors, where it remains for years.
\end{description}

\caution{Flash memory (\enquote{flash}) has nothing to do with a pipeline
flush (Lesson 3).}

A flash SSD is slower than DRAM but still much faster than a magnetic disk,
needs little power, and retains its data without a battery.\margin{Flash
cannot be overwritten in place: a block of several MB must be erased first,
and survives only some thousand erases.}  Its cost per
gigabyte lies between those of DRAM and disk, so for the same price it
offers much less capacity than a disk: for a long time, typical flash SSDs
held gigabytes where disks held terabytes.

\section{Hybrid magnetic--flash disks}
\label{sec:l16-hybrid}

Magnetic disks and flash memory have complementary
strengths.\source{Lesson 16, videos 11--12.}  A magnetic disk offers a huge
capacity at a low cost per gigabyte, but it is slow, consumes considerable
power while its platters spin, and is vulnerable to shocks during operation.
Flash memory is fast, consumes little power and has no moving parts, but it
costs more per gigabyte.

\begin{definition}[Hybrid disk]
  A \term{hybrid disk}[ハイブリッドディスク] combines a magnetic disk with a
  smaller amount of flash memory in one device.  The magnetic disk holds all
  data, and the flash memory serves as a cache for the data that is used most
  frequently (\cref{fig:l16-hybrid}).
\end{definition}

Temporal and spatial locality (Lesson 12) apply to files as they do to memory
accesses: the operating system, the applications in daily use and the files
being worked on account for most accesses.  If they are kept in flash, most
requests are served at the speed of flash and without mechanical movement,
while the capacity and the cost per gigabyte remain close to those of the
disk.  When requests are served from flash, the disk can even stop spinning,
which saves power and makes the device less sensitive to shocks; a request
that then does need the disk waits, often for seconds, until the platters
spin again.

\begin{marginfigure}
  % Hybrid disk: a controller that serves most requests from a flash cache
% holding the frequently used data; the magnetic disk holds all data.
\resizebox{\linewidth}{!}{%
\begin{tikzpicture}[x=1cm,y=1cm, font=\footnotesize\sffamily, >={Stealth[length=1.6mm]},
  box/.style={draw, fill=ln-box, align=center, inner sep=3pt, minimum height=0.9cm},
  lab/.style={font=\scriptsize\sffamily, text=ln-muted, align=center}]
  \node[box, minimum width=4.2cm] (ctl) at (0,0) {\iflangja{コントローラ}{controller}};
  \node[box, fill=ln-accent!20, minimum width=1.9cm] (flash) at (-1.15,-1.9)
    {\iflangja{フラッシュ}{flash}\\[-1pt]{\scriptsize\iflangja{よく使うデータ}{frequently used}}};
  \node[box, minimum width=1.9cm] (disk) at (1.15,-1.9)
    {\iflangja{磁気ディスク}{magnetic disk}\\[-1pt]{\scriptsize\iflangja{全データ}{all data}}};
  \draw[<->, very thick, ln-accent] (-1.15,-0.45) -- (flash.north)
    node[lab, midway, left] {\iflangja{大部分}{most}};
  \draw[<->] (1.15,-0.45) -- (disk.north)
    node[lab, midway, right] {\iflangja{残り}{rest}};
  \draw[<->, thick] (0,1.2) -- (ctl.north) node[lab, pos=0.25, right] {\iflangja{I/O バス}{I/O bus}};
  \node[draw=ln-muted, dashed, rounded corners=2pt, fit=(ctl)(flash)(disk), inner sep=5pt] (hyb) {};
  \node[lab, anchor=north] at (hyb.south) {\iflangja{ハイブリッドディスク}{hybrid disk}};
\end{tikzpicture}}

  \caption{A hybrid disk.  Most requests are served by the flash cache.}
  \label{fig:l16-hybrid}
\end{marginfigure}

The average access time behaves like the AMAT of a cache
(\cref{eq:l12-amat}).  A request served from flash takes the hit time, and a
request served by the disk takes the total time of a miss.  The controller
knows where each block lies, so a miss does not read the flash first, and the
miss penalty is the difference between the two times.

\begin{example}
  \label{ex:l16-hybrid}
  A hybrid disk combines flash with an access time of
  \SI{0.1}{\milli\second} and the disk of \cref{ex:l16-access}, which is kept
  spinning.  For the request of that example, a hit also pays the controller
  and I/O bus times, $0.1 + 0.3 + 0.5 = \SI{0.9}{\milli\second}$, and a miss
  takes \SI{9}{\milli\second}, so the miss penalty is
  \SI{8.1}{\milli\second}.  If \SI{95}{\percent} of the requests find their
  data in flash, the average access time is $0.9 + 0.05 \times 8.1 \approx
  \SI{1.3}{\milli\second}$, about 7 times shorter than with the disk alone,
  although the flash holds only a small part of the data.  The controller and
  the I/O bus now account for most of the time of a hit.\tip{The hit time is
  a floor: however high the hit rate, this disk cannot beat $9/0.9 = 10$
  times the disk alone (Amdahl's law, Lesson 2).}
\end{example}

\Cref{tab:l16-compare} summarizes the technologies of this and the previous
section.  Which one is preferable depends on the requirement that dominates:
the lowest cost for a large amount of data favors a magnetic disk, speed,
low power and robustness favor flash, and a hybrid disk approaches both for
workloads whose accesses concentrate on a small part of the data.\caution{The
flash is a cache, not extra capacity: a 2~TB hybrid disk with 8~GB of flash
is sold, and behaves, as a 2~TB drive.}

\begin{table}[htbp]
  \centering\small
  \caption{Magnetic disks and semiconductor storage compared.}
  \label{tab:l16-compare}
  \begin{tabular}{@{}lllll@{}}
    \toprule
                        & Magnetic disk & DRAM SSD     & Flash SSD & Hybrid disk \\
    \midrule
    Access              & slow (ms)     & fastest      & fast      & fast for most \\
    Cost per GB         & low           & very high    & medium    & low \\
    Power               & high          & high         & low       & lower \\
    Moving parts        & yes           & no           & no        & yes \\
    Keeps data unpowered & yes          & while battery lasts & yes & yes \\
    \bottomrule
  \end{tabular}
\end{table}

\needspace{6\baselineskip}
\begin{keypoint}
  A hybrid disk applies the principle of caching to storage: flash in front
  of a magnetic disk gives most requests the speed and low power of flash at
  nearly the cost per gigabyte of the disk.
\end{keypoint}

\section{Connecting I/O devices}
\label{sec:l16-io}

Storage devices, like other input/output (I/O) devices, are not attached
directly to the processor.\source{Lesson 16, videos 13--14.}  They are
connected through buses defined by industry standards.

\begin{definition}[I/O bus]
  An \term{I/O bus}[I/O バス] is a standardized bus through which I/O
  devices are connected to a computer.  The standard fixes the connectors,
  the electrical signals and the protocol; PCI Express, SATA, SCSI and USB
  are examples.
\end{definition}

Standardization lets any device built to the standard work with any computer
that provides the bus, whoever manufactured either of them, and a computer
maker can offer one kind of interface for a wide range of devices.  The
price is a limited rate of improvement.  A new version of a standard has to
be agreed on by many companies, and devices built for the older version
should keep working, so a bus cannot adopt a faster technology as soon as it
becomes available.  The more devices depend on a standard, the more slowly
it changes.\sidenote{PCI Express doubles its per-lane rate each
generation: \num{2.5}~GT/s in 2003, then 5, 8, 16 (a seven-year wait), 32
and 64~GT/s in 2022.  Compatibility runs both ways: a new card works in an
old slot, at the old rate.}

I/O devices also differ enormously in the throughput they need: a graphics
adapter moves huge amounts of data, a disk much less, a keyboard almost none.
A single fast bus would force every device to have an expensive fast
interface, and a single slow bus would throttle the fast devices.  The buses
are therefore arranged in a hierarchy (\cref{fig:l16-io-buses}).

\begin{definition}[Mezzanine bus]
  A \term{mezzanine bus}[メザニンバス] is a fast, short I/O bus between the
  processor and the slower I/O buses; today it is PCI Express.  Fast devices
  such as the graphics adapter are connected to it directly, and the
  controllers of the slower buses are attached to it.
\end{definition}

A SATA or SCSI controller on the mezzanine bus provides a SATA (or SCSI) bus
for disks and optical drives, and a USB host controller provides a USB
bus,\margin{SATA~3 carries 6~Gb/s (600~MB/s), USB 3.2 Gen~2 10~Gb/s and a
16-lane PCIe~5.0 slot some 63~GB/s; a keyboard needs a few hundred bytes per
second.}
extended by hubs, for external devices such as keyboards and flash drives.
These buses serve naturally slower devices, for which a stable standard
matters more than speed; they are slower than the mezzanine bus and change
less often.  In PC chipsets the SATA and USB controllers are integrated into
one chip, the \term{southbridge}[サウスブリッジ], which is connected to the
northbridge (Lesson 15).

\begin{figure}[htbp]
  \centering
  % Hierarchy of buses (logical view): processor and northbridge with main
% memory; the mezzanine bus (PCI Express) connects the graphics adapter
% directly and the SATA and USB controllers, which provide the slower buses
% (two point-to-point SATA links; a USB bus fanning out to its devices).
% In PC chipsets the two controllers form the southbridge (dashed box).
\begin{tikzpicture}[x=1cm,y=1cm, font=\footnotesize\sffamily, >={Stealth[length=1.6mm]},
  box/.style={draw, fill=ln-box, align=center, inner sep=3pt, minimum height=0.7cm},
  ctl/.style={draw, fill=ln-box, align=center, inner sep=2pt, minimum height=0.6cm, minimum width=2.1cm, font=\scriptsize\sffamily},
  dev/.style={draw, fill=white, align=center, inner sep=2pt, minimum height=0.6cm, font=\scriptsize\sffamily},
  lab/.style={font=\scriptsize\sffamily, text=ln-muted}]
  \node[box, minimum width=2.4cm] (cpu) at (5.0,3.3) {\iflangja{プロセッサ}{processor}};
  \node[box, minimum width=2.4cm, fill=ln-accent!20] (nb) at (5.0,2.05) {\iflangja{ノースブリッジ}{northbridge}};
  \node[dev, minimum width=1.9cm] (mem) at (0.95,2.05) {\iflangja{メインメモリ}{main memory}};
  \draw[<->, very thick, ln-accent] (cpu) -- node[lab, right] {FSB} (nb);
  \draw[<->, very thick, ln-accent] (nb) -- node[lab, above] {\iflangja{メモリバス}{memory bus}} (mem);
  % mezzanine bus
  \coordinate (bus) at (0,1.0);
  \draw[very thick, ln-accent] (nb.south) -- (nb.south |- bus);
  \draw[very thick, ln-accent] (1.25,1.0) -- (8.0,1.0);
  \node[lab, anchor=south west] at (5.12,1.02) {\iflangja{メザニンバス（PCI Express）}{mezzanine bus (PCI Express)}};
  \node[dev, minimum width=1.9cm] (gfx) at (1.25,0.0) {\iflangja{グラフィックス}{graphics}};
  \node[ctl] (sata) at (4.35,0.0) {\iflangja{SATA コントローラ}{SATA controller}};
  \node[ctl] (usb) at (8.0,0.0) {\iflangja{USB コントローラ}{USB controller}};
  \foreach \n in {gfx,sata,usb} {
    \draw[very thick, ln-accent] (\n.north) -- (\n.north |- bus);
  }
  % southbridge: the controllers of the slower buses
  \node[draw=ln-muted, dashed, rounded corners=2pt, fit=(sata)(usb), inner sep=4pt] (sb) {};
  \node[lab, anchor=west] at (sb.east) {\iflangja{サウスブリッジ}{southbridge}};
  % slower buses: two point-to-point SATA links, one USB bus
  \node[dev, minimum width=1.6cm] (d1) at (3.3,-1.55) {\iflangja{磁気ディスク}{magnetic disk}};
  \node[dev, minimum width=1.6cm] (d2) at (5.4,-1.55) {\iflangja{光ディスク装置}{optical drive}};
  \node[dev, minimum width=1.3cm] (d3) at (7.25,-1.55) {\iflangja{キーボード}{keyboard}};
  \node[dev, minimum width=1.3cm] (d4) at (8.75,-1.55) {\iflangja{USB メモリ}{flash drive}};
  \draw[thick] ([xshift=-3mm]sata.south) -- node[lab, left] {SATA} (d1.north);
  \draw[thick] ([xshift=3mm]sata.south) -- node[lab, right] {SATA} (d2.north);
  \coordinate (uj) at (8.0,-0.8);
  \draw[thick] (usb.south) -- (uj);
  \draw[thick] (uj) -- (d3.north);
  \draw[thick] (uj) -- (d4.north);
  \fill (uj) circle[radius=0.05];
  \node[lab, anchor=west] at (8.08,-0.6) {USB};
\end{tikzpicture}

  \caption{The hierarchy of buses.  The graphics adapter is connected
    directly to the mezzanine bus; disks and other slower devices are
    reached through the SATA and USB controllers on it, which provide
    slower standard I/O buses.  In PC chipsets these controllers form the
    southbridge (dashed).}
  \label{fig:l16-io-buses}
\end{figure}

Each controller translates between the bus it provides and the mezzanine
bus.  The data of a disk read travels from the disk over SATA to the SATA
controller, over the mezzanine bus and through the northbridge into main
memory; this path is the I/O bus time of \cref{eq:l16-access}.  Fast devices
sit close to the processor, and the many slow devices are confined to their
own buses behind the controllers, where they do not slow down the fast
ones.\margin{Newer processors integrate the northbridge and connect to the
southbridge directly.}

Whatever the technology and however it is connected, storage is where data
is expected to survive, yet every storage device eventually fails.  Lesson
17 shows how a system can keep its data and its service despite such
failures.

\begin{keypoint}[Lesson summary]
  Storage holds files and the pages of virtual memory persistently, and its
  latency improves even more slowly than that of DRAM.  A magnetic disk
  stores data in sectors on the tracks of rotating platters
  (\cref{eq:l16-capacity}); an access (\cref{eq:l16-access}) is dominated by
  the seek time and the rotational latency, half a revolution on average,
  and this mechanical latency improves only slowly, while capacity grew
  quickly for many years.  Optical disks are portable but improve slowly
  because of standardization; magnetic tape offers cheap capacity with
  sequential access
  and serves for backup.  Solid-state disks built from DRAM or flash memory
  avoid mechanical movement at a higher cost per gigabyte, and a hybrid disk
  uses flash as a cache for a magnetic disk.  I/O devices are connected
  through standardized I/O buses arranged in a hierarchy: fast devices sit
  directly on a fast, short mezzanine bus such as PCI Express, and slower
  devices on slower, more stable buses such as SATA and USB, whose
  controllers are attached to the mezzanine bus.
\end{keypoint}

\end{document}
